\section{Demonstration: three generator topologies under identical winding faults}\label{sec:demo}

\subsection{Alternatives and data}\label{sec:data}

The three alternatives are laboratory-scale generators of the same power class on instrumented test benches at the University of S\~ao Paulo, whose records were published as open benchmark datasets (Table~\ref{tab:alternatives}). The PMSG \citep{Tominaga2025pmsg} and the SCIG \citep{Tominaga2025scig} are four-pole, 230\,V, double-star machines from the same manufacturer, rated 2.5\,kVA and 2.5\,kW respectively, installed on the \emph{same} bench: the same induction-motor prime mover and frequency converter, the same rapid-control-prototyping platform and field-oriented control (rotor angle from an encoder for the PMSG, from a flux observer for the SCIG), the same sensors, the same acquisition (20\,kHz, 3\,s records), the same fault-insertion hardware and the same test script. The variable-speed WFSG \citep{Tominaga2024gen} is a 2\,kVA four-salient-pole machine on an earlier bench of the same group \citep{Rocha2023} with a three-level converter, field-oriented control with encoder angle, 4\,kHz acquisition and 1.3\,s records. The PMSG--SCIG pair is therefore a controlled comparison in which everything but the machine and the control structure it requires is held fixed; the WFSG is a third alternative with the same fault protocol but a different bench, and is analysed with that difference kept explicit. Terminal voltages are measured at the machine terminals by each bench's voltage transducers, whose bandwidth the data papers do not document, so the terminal-voltage features carry the sensing chain of their bench as well as the machine.

\begin{table}[t]
\footnotesize
\setlength{\tabcolsep}{4pt}
\caption{The three alternatives. PMSG and SCIG share one bench and one protocol; the WFSG bench differs where indicated. Ratings as given by the manufacturer (kVA for the synchronous machines, kW for the induction machine).}\label{tab:alternatives}
\begin{tabular}{@{}p{2.4cm}p{2.8cm}p{3.3cm}p{3.4cm}@{}}
\toprule
 & PMSG & SCIG & WFSG (variable speed) \\
\midrule
Rating, poles & 2.5\,kVA, 4 & 2.5\,kW, 4 & 2\,kVA, 4 salient \\
Excitation & permanent magnets & induction (magnetising current) & wound field, static exciter \\
Converter and control & 2L, FOC with encoder & 2L, FOC with flux observer & 3L-NPC, FOC with encoder \\
Acquisition & 20\,kHz, 3\,s & 20\,kHz, 3\,s & 4\,kHz, 1.3\,s \\
Reference / faulted interval & 0.68\,s / 0.36\,s & 0.68\,s / 0.36\,s & 0.49\,s / 0.06--0.28\,s \\
Fault impedance & 2.6\,$\Omega$ & 2.6\,$\Omega$ & 1, 2.83, 5.66, 11.32\,$\Omega$ \\
Operating points & 3 speeds $\times$ 3 torques & same & 716--1910\,rpm, 4 torque levels (sparse) \\
Healthy trials & 9 & 9 & none \\
Fault trials & 216 & 216 & 438 (198 at 2.83\,$\Omega$) \\
\botrule
\end{tabular}
\end{table}

\subsection{Failure-mode class and extent scale}\label{sec:extent}

All three machines were built with 24 tapped derivation points along their stator windings at the same relative positions, so that a short-circuit between two taps realises a fault of known extent (Online Resource~1, Table~S9); the tap positions, given as percentages of a winding branch from the neutral, are identical in the three data papers, which is what makes the extent scale realisable in every alternative. Two classes are used: inter-turn short-circuits (both taps on the same branch) and inter-winding short-circuits (taps on the two parallel branches of a phase). The extent $e$ is the fraction of a branch between the two taps; it ranges from 2.7\,\% to 12.1\,\% for the twelve inter-turn cases and from 2.3\,\% to 40.2\,\% for the twelve inter-winding cases. Extents that coincide at different winding locations (2.8\,\% twice, 7.4\,\% three times, 11.6\,\% twice) form one extent level. Faults were closed through a contactor and a series resistance after the machine had reached steady state: 2.6\,$\Omega$ for 0.4\,s on the PMSG and SCIG; on the WFSG the experimenters raised the resistance with extent (1\,$\Omega$ for the inter-turn cases below 11\,\%, 2.83\,$\Omega$ for the remaining inter-turn cases and the small inter-winding cases, 5.66 and 11.32\,$\Omega$ for the large inter-winding cases), so that on that bench extent and impedance are confounded. The fault current is measured on all benches; it is used only as ground truth for the physical severity axis, never as a feature. Because the series resistance is of the order of the winding impedance of these small machines, the fault current is essentially the voltage induced in the shorted turns divided by the resistance, a \emph{resistor-limited} severity that differs from a bolted fault, which is limited by the leakage reactance of the shorted turns (Section~\ref{sec:res-severity}).

The PMSG and SCIG matrix is complete: 24 fault cases $\times$ 3 speeds (1200, 1500, 1800\,rpm) $\times$ 3 torques (5.2, 6.4, 8.0\,N\,m), the torque--speed pairs scaled from the operating region of the IEA 15\,MW reference turbine \citep{Gaertner2020} to the laboratory speed range (the electrical frequency is not scaled), plus the healthy condition at each operating point; the SCIG generates with slip, at 59.6\,Hz at 1800\,rpm. The WFSG matrix covers the same 24 tap pairs at six speeds (716--1910\,rpm) and four per-unit torque levels, sparsely and with two to six repetitions, so its operating points are not matched to the other two. The WFSG benchmark also contains phase-to-phase faults, a class absent from the PMSG/SCIG protocol; they are excluded from every analysis. The analyses use the WFSG's 438 inter-turn and inter-winding trials at all impedances for the null, the feature ranking and the detectable shares, with the impedance as a covariate, and the 198 trials at 2.83\,$\Omega$---the impedance closest to the PMSG/SCIG value, covering four inter-turn tap pairs of 11.6--12.1\,\% and four inter-winding pairs of 2.3--10.6\,\%---for the minimum detectable extent and the decision table. Each table and figure states which set it uses.

\subsection{Observation suites and features}\label{sec:suites}

Five observation suites are considered, ordered by the data access, hardware and integration they add to a bare generator--converter set (Table~\ref{tab:suitedef}). \emph{Drive-internal} uses only quantities the converter controller already has: the $dq$ currents, its own $dq$ voltage commands, the proportional--integral controller outputs, the modulation duty cycles, the DC-link voltage, the estimated electrical torque and the speed. \emph{3\,CT} uses the three phase currents---from dedicated transducers or from the converter's own current sensors routed to a monitoring system---with angle-free features only; for a converter-fed machine its cost is data access and bandwidth rather than new hardware. \emph{3\,CT + 3\,VT} adds three terminal voltage transducers; \emph{3\,CT + 3\,VT + drive} combines the external transducers with controller access, which also allows the measured terminal voltages to be transformed into the synchronous frame; and \emph{+ torque transducer} adds a shaft torque sensor. For the WFSG the field current is an additional drive-internal observable (the \emph{excitation} suite), available because its exciter is static; a brushless exciter would not expose it \citep{Neti2009}. Each feature is assigned to the cheapest suite whose measurements suffice to compute it. The suites are those the datasets can populate: zero-sequence voltage (which needs neutral access), stray flux, winding temperature, vibration and partial discharge are not recorded, and the angle-free Park's-vector modulation of the currents was not computed in this version; the suite conclusions therefore hold among the suites and features considered. The zero-sequence current is not used: the machines have isolated star points, so it is a noise floor.

\begin{table}[t]
\footnotesize
\setlength{\tabcolsep}{4pt}
\caption{Observation suites and the features they can compute. $f_e$: electrical frequency; ``2$f_e$'': amplitude of the component at twice the electrical frequency; $I_1, I_2$ ($V_1, V_2$): positive- and negative-sequence magnitudes at $f_e$.}\label{tab:suitedef}
\begin{tabular}{@{}p{2.3cm}p{3.0cm}p{7.0cm}@{}}
\toprule
Suite & Added access / hardware & Features \\
\midrule
Drive-internal & controller data & $I_d, I_q$: std, 1$f_e$, 2$f_e$; $V_d^{\mathrm{cmd}}, V_q^{\mathrm{cmd}}$: 2$f_e$; PI outputs $d,q$: std, 2$f_e$; duty-cycle negative sequence $D_2/D_1$; $T_e$: std, 1$f_e$, 2$f_e$; DC-link std; speed std \\
3\,CT & three phase currents, no angle & $I_2/I_1$, RMS unbalance, harmonics 3, 5, 7 of $I_a$ relative to the fundamental, THD \\
3\,CT + 3\,VT & + 3 voltage transducers & 3\,CT features + $V_2/V_1$ \\
3\,CT + 3\,VT + drive & + controller data & all of the above + measured $V_d, V_q$: 2$f_e$ \\
+ torque transducer & + shaft torque sensor & + $T_m$: std, 1$f_e$, 2$f_e$ \\
Excitation (WFSG) & field-current measurement & $I_f$: std, 2$f_e$ \\
\botrule
\end{tabular}
\end{table}

Features are computed on sliding windows of $n_c$ electrical cycles with a hop of one cycle. The electrical frequency and the phase rotation are estimated per window from the phase slope of the current space vector (falling back to the nominal value outside 0.5--1.5 times it), so that no angle sensor is needed for the external suites; sequence components are obtained by Hann-windowed projection of each phase quantity onto the fundamental and symmetrical-component transformation; 1$f_e$ and 2$f_e$ amplitudes of synchronous-frame quantities by projection after mean removal. On the WFSG bench the measured converter-side phase voltages serve as terminal voltages, the synchronous-frame voltages are obtained from them with the encoder angle, the commanded voltages are the controller's, and the electrical torque is computed from power and speed. Online Resource~1 (Section~S1) lists the definitions.

\subsection{Computation}\label{sec:computation}

For the PMSG and SCIG the reference interval is the 0.68\,s before the contactor command, split into $R_1$ (0.32\,s) and $R_2$ (0.36\,s), and the faulted interval runs from 30\,ms after the command to 10\,ms before its release; because the relay picks up 85--90\,ms after the command, that interval contains about 55\,ms of healthy time, which dilutes the fault-interval mean and makes the PMSG/SCIG SDRs slightly conservative. For the WFSG the fault interval is located from the measured fault current and the reference interval, 0.49\,s, is split at 0.30\,s. The main PMSG--SCIG comparison uses $n_c=5$; the three-alternative comparison uses $n_c=3$ for all three machines because the WFSG fault intervals are short: at 23.9--63.7\,Hz a faulted interval holds between zero and eleven three-cycle windows, four trials hold none and are dropped, and 81 of 438 hold fewer than three. The sensitivity of every conclusion to $n_c\in\{3,5,8\}$ and to $\alpha\in\{0.90,0.95,0.99\}$ is reported, with the screen applied at each setting. The null is pooled over all 225 records of each of the PMSG and SCIG and over the 438 WFSG records, every record contributing a healthy reference interval. Uncertainty is quantified by bootstrap \citep{Efron1993} with 2000 resamples: the median SDR and the detectable share of a class are resampled over the twelve tap pairs, the same pairs for both machines so that the SCIG/PMSG ratio is matched; the minimum detectable extent is resampled over the operating-point trials within each extent level; shares also carry Wilson intervals; the null quantile is held at its full-sample value. The decision table uses 1000 resamples per suite, alternative, class and $\alpha$, with $e_{\mathrm{req}}\in\{3,5,8,12\}$\,\%. All interval boundaries, thresholds, seeds and model settings are listed in Online Resource~1 (Section~S2).

Step 6 of the procedure is instantiated as a window-level change detector: a gradient-boosted tree classifier \citep{Friedman2001} and, for the PMSG--SCIG pair, a logistic regression, trained to separate faulted windows from healthy ones. Within an alternative the detector is evaluated by five-fold cross-validation grouped by tap pair (healthy trials grouped by speed, so that every fold contains some), so that the tested fault locations are never seen in training; across alternatives it is trained on all trials of one machine and tested on all trials of another. Three feature-referencing rules are compared: physical units; features $z$-scored on the target machine's healthy trials (a commissioning baseline); and the absolute $z$-score of each window against the mean and standard deviation of the \emph{first half} of its own trial's reference interval. Negatives are the reference and recovery windows of all trials plus the fault-interval windows of the healthy trials, the only negatives that contain the contactor operation without a fault. Because the windows used to normalise a trial are trivially separable under the third rule, every metric is reported on all windows and on the out-of-reference windows, and the false-alarm rate of the threshold that gives 1\,\% false positives on all negatives is reported on the healthy trials' fault-time windows. The true-positive rate at 1\,\% false-positive rate is a point of the receiver operating characteristic, not an operating threshold.

All computations are implemented in Python and released with the paper; the pipeline runs from the feature caches to every table and figure in under thirty minutes on a four-core machine. A large language model was used as an assistant in writing the analysis code and in drafting and editing the text; every analysis and every reference was verified by the author (see the declarations).
